\section{Introduction}

A succinct non-interactive argument of knowledge (SNARK) allows a
prover to establish knowledge of a witness using a short proof that
can be verified much faster than the computation it
certifies~\cite{innovations/BitanskyCCT12}.
SNARKs enable efficient verification of computations performed by
untrusted parties; zero-knowledge variants also hide the
witness~\cite{eurocrypt/Groth16,crypto/Setty20}.
The cost of producing these proofs, however, can substantially
exceed the cost of the computation itself.

For computations on short values, the proof's arithmetic field is
one source of this overhead. Many efficient SNARKs combine a
polynomial interactive oracle proof (PIOP) with a polynomial
commitment scheme (PCS)~\cite{eurocrypt/BunzFS20,eurocrypt/ChiesaHMMVW20}.
The PIOP expresses correctness through polynomial identities,
and the PCS authenticates the evaluations
used to check them~\cite{crypto/Setty20,eurocrypt/ChenBBZ23}.
Checking these identities with small soundness error requires a
challenge field large enough relative to their degrees.
Its elements can therefore be much larger than the witness entries.

The effect extends beyond storing or committing the witness.
In sumcheck, random challenges combine witness-table entries into
values that generally occupy the full challenge field, even when
the original entries lie in a small subfield~\cite{iacr/BagadDDT25}.
Subsequent rounds multiply and store these larger values. Thus a
small-field witness can still incur large-field arithmetic and
memory costs throughout much of the PIOP.

Quantized machine-learning inference makes this disparity concrete.
SmoothQuant supports matrix multiplications with 8-bit weights and
8-bit activations~\cite{icml/XiaoLSWDH23}.
In zero-knowledge machine learning (zkML), Kang et al.~\cite{corr/KangHSS22}
prove DNN inference with 8-bit weights and activations and
intermediate computations of up to 32 bits over a 254-bit prime field.
The field used to prove the inference is much larger than the
values on which the model operates.

Field-agnostic commitments permit the arithmetic field to match
the application~\cite{crypto/GolovnevLSTW23,crypto/ZeilbergerCF24},
and small-field commitments can retain the
witness's compact representation. The challenge-dependent
computations in the PIOP still require a way to exploit that
representation~\cite{eurocrypt/DiamondP26,iacr/BagadDDT25}.
We therefore ask:
\begin{quote}
\emph{Can a SNARK retain small-field arithmetic for most of the
prover's work while drawing its soundness challenges from a large field?}
\end{quote}

\paragraph{Our approach.}
We present \methodname, a framework that uses \emph{polynomial
packing} to reduce the number of entries processed by the
large-field proof. Let \(N=m\kappa\) and
\([\ExtensionField:\Field]=\kappa\).
Our central reduction transforms a length-\(N\) vector relation
over the small field \(\Field\) into a length-\(m\) relation over
\(\ExtensionField\).
We arrange the witness as an \(m\times\kappa\) matrix and encode
each row by a polynomial over \(\Field\). A low-degree identity
captures the original relation, and its evaluation at a random
extension-field point gives the shorter relation.
The prover computes the row-polynomial products before this
evaluation, keeping these multiplications over \(\Field\).
Only the resulting length-\(m\) relation is proved over
\(\ExtensionField\); Section~\ref{sec:overview} gives the details.

We construct such encodings for inner and Hadamard products,
reducing both to a length-\(m\) extension-field inner product,
and use the Hadamard reduction to prove permutation constraints.
For \(|\Field|\geq\kappa\) and \(m\geq\kappa\), both reductions
use \(O(\kappa N)\) base-field operations and \(O(N)\) base-field
coordinates of memory under schoolbook arithmetic and precomputed
bases, excluding commitments and openings.
The corresponding uncompressed extension-field proof uses
\(O(\kappa^2N)\) base-field operations and \(O(\kappa N)\)
coordinates.

To authenticate the shorter relation, we must connect its
challenge-dependent packing to a commitment fixed before the
challenge. Our packed PCS commits using a fixed field basis and
verifies the resulting evaluation claims through base-field
coordinate identities.
We compose the product, permutation, and opening arguments into
protocols for a single circuit execution and for multiple
executions of the same circuit. In the latter case, packing
corresponding wires preserves the shared selectors and wiring.
Theorems~\ref{theorem:single-instance-packing}
and~\ref{theorem:multiple-instance-packing} summarize their
efficiency and soundness.

The reductions also admit binary arithmetizations. The tower
\(\mathbb F_2\subset\mathbb K=\mathbb F_{2^8}\subset\ExtensionField\)
allows byte-field arithmetic in the PIOP while retaining a larger
challenge field.
Section~\ref{sec:binary-packing} gives the required parameters and
Booleanity constraints, and explains how to use a compatible
code-based binary PCS through the code-switching interface of
Section~\ref{sec:pcs}.

Our component benchmarks compare the packed reductions with a
small-field-first baseline. At \(N=2^{30}\) over an 8-bit base
field, the measured prover speedups are \(5.60\times\) for inner
products and \(4.31\times\) for Hadamard products, with a
\(4.48\times\) reduction in resident memory for both.
These measurements include packing and witness-dependent
preprocessing, and exclude commitments and PCS openings.
Section~\ref{sec:experiment} reports the full comparison across
8-, 16-, 32-, and 64-bit configurations and the measurement procedure.

\subsection{Related Work}
\label{sec:related-work}

\paragraph{SNARKs and polynomial IOPs.}
Our construction follows the separation of algebraic verification
and cryptographic authentication provided by the PIOP--PCS
framework. This framework builds on interactive oracle
proofs~\cite{tcc/Ben-SassonCS16} and their polynomial-oracle
formulation~\cite{eurocrypt/BunzFS20,eurocrypt/ChiesaHMMVW20}.
The prover commits to its oracle polynomials and authenticates
the verifier's queries with PCS evaluation proofs.
For suitable public-coin instantiations, the
Fiat--Shamir transform~\cite{crypto/FiatS86} removes interaction in
the random-oracle model~\cite{tcc/Ben-SassonCS16,eurocrypt/BunzFS20}.
Spartan~\cite{crypto/Setty20} uses sumcheck to verify rank-one
constraint systems, while HyperPlonk~\cite{eurocrypt/ChenBBZ23}
expresses PLONK-style gate and permutation constraints over the
Boolean hypercube. Both use multilinear polynomial representations.
These multilinear interfaces allow the relation proof and its
commitment scheme to be optimized separately.

\paragraph{ECC-based constructions.}
Group-based constructions offer short proofs with efficient
verification, but tie the arithmetic field to the chosen group.
Quadratic span programs and quadratic arithmetic
programs~\cite{eurocrypt/GennaroGP013} provide polynomial
characterizations of circuit satisfiability, underpinning early
preprocessing zkSNARKs such as Pinocchio~\cite{cacm/ParnoHG016}.
Groth's pairing-based SNARK~\cite{eurocrypt/Groth16} has proofs
of three group elements and circuit-specific preprocessing.
Polynomial commitments provide a modular route to such arguments.
KZG~\cite{asiacrypt/KateZG10} introduced constant-size commitments
and evaluation proofs for univariate polynomials, and Papamanthou,
Shi, and Tamassia~\cite{tcc/PapamanthouST13} extended the primitive
to multivariate polynomials.
Sonic~\cite{ccs/MallerBKM19}, Marlin~\cite{eurocrypt/ChiesaHMMVW20},
and PLONK~\cite{iacr/GabizonWC19} provide SNARKs with universal,
updatable structured reference strings.
In the discrete-logarithm setting,
Bulletproofs~\cite{sp/BunzBBPWM18} gives an inner-product argument
with logarithmic communication and no trusted setup, and
Hyrax~\cite{sp/WahbyTSTW18} provides a multilinear commitment with
communication--verification tradeoffs.
Their polynomial arithmetic takes place in the group's scalar
field, whose size is constrained by cryptographic security
independently of the witness bit width.

\paragraph{Hash-based constructions.}
Hash-based arguments instead authenticate error-correcting encodings
with Merkle trees, separating the arithmetic field from a group
order. STARKs~\cite{iacr/Ben-SassonBHR18} use this approach
with FRI~\cite{icalp/Ben-SassonBHR18}, an interactive oracle proof
of proximity for Reed--Solomon codes.
Ligero~\cite{dcc/AmesHIV23} uses interleaved Reed--Solomon codes
to obtain sublinear arguments.
Aurora~\cite{eurocrypt/Ben-SassonCRSVW19} gives a transparent
IOP-based zkSNARK for R1CS, while
Fractal~\cite{eurocrypt/ChiesaOS20} uses holographic IOPs to
support transparent preprocessing and recursive composition.
WHIR~\cite{eurocrypt/ArnonCFY25} improves verification through
constrained Reed--Solomon proximity tests that support evaluation
queries. Efficient Reed--Solomon encoding can nevertheless impose
algebraic conditions on the field.

\paragraph{Field-agnostic commitments and SNARKs.}
Field-agnostic constructions remove this restriction: they work
over any sufficiently large finite field, without a prescribed
group order or an FFT-friendly multiplicative subgroup.
Brakedown~\cite{crypto/GolovnevLSTW23} obtains linear-time proving
from linear-time encodable codes.
Block et al.~\cite{crypto/BlockFKTWZ24} generalize
expand-accumulate codes to arbitrary finite fields and instantiate
Brakedown's commitment framework with these codes.
BaseFold~\cite{crypto/ZeilbergerCF24} combines foldable codes with
sumcheck to obtain a field-agnostic multilinear PCS with
polylogarithmic verification.
Code switching~\cite{jacm/RonZewiR24} separates the choice of an
efficient encoding from the mechanism used to verify evaluations.
Blaze~\cite{eurocrypt/BrehmCFRRZ25} uses this approach with
interleaved RAA codes to obtain an efficient PCS over binary fields.

These developments broaden the choice of arithmetic field and
improve the commitment layer. Once the PIOP draws challenges from
a larger field, however, its intermediate values can still occupy
that field. Retaining the benefit of short witness entries also
requires exploiting their structure within the relation
proof~\cite{eurocrypt/DiamondP25,iacr/BagadDDT25}.

\paragraph{Small-field embeddings and Binius.}
Field embeddings offer another way to exploit small witness values.
Cascudo and Giunta~\cite{iacr/CascudoG21} use reverse
multiplication-friendly embeddings to verify Boolean R1CS through
larger-field linear and row checks.
Binius, due to Diamond and Posen~\cite{eurocrypt/DiamondP25},
adapts Brakedown to multilinear commitments over binary towers
without coefficient-wise embedding overhead, and gives corresponding
product, permutation, and lookup arguments.
Their subsequent ring-switching construction~\cite{eurocrypt/DiamondP26}
reduces multilinear commitments over a ground field to commitments
over an extension field. Its commitment cost matches that of the
underlying PCS on inputs of the same bit size, and its opening
reduction adds linear prover and logarithmic verifier overhead.
The compiler applies to general field extensions, with a concrete
PCS based on a binary-field adaptation of BaseFold.
Ring-switching supplies the packed commitment interface; Binius's
relation-specific arguments allow a circuit PIOP to use the same
compact representation.

\paragraph{Small-value and small-field sumcheck.}
The cost of the relation proof is also the subject of work on
sumcheck over short values.
Bagad, Domb, and Thaler~\cite{iacr/BagadDT24} optimize sumcheck
for input polynomials whose values lie in a small subfield, replacing
extension-field multiplications with additional base-field arithmetic.
Bagad, Dao, Domb, and Thaler~\cite{iacr/BagadDDT25} extend this
approach to short integers and small-subfield values, combine it
with equality-polynomial optimizations, and integrate the prover
into Jolt.
Dao and Thaler~\cite{iacr/DaoT24} exploit binary towers in
zero-check, reducing equality-table computation and generalizing
constraint packing with linearly independent extension-field weights.
Wei et al.~\cite{iacr/WeiWXZDWZ25} give packed sumcheck over small
prime fields using repetition and folding to retain base-field
computation, with an application to TFHE bootstrapping.
This line of work exploits the input values and the equality
polynomial to reduce the cost of sumcheck proving.

\paragraph{Relation to our work.}
\methodname reduces the dimension of the relation before invoking
the extension-field proof. For inner and Hadamard products, we
encode the relation in a polynomial identity whose row products
are computed over the base field. Random evaluation then reduces
the original length-\(N\) claim to a length-\(N/\kappa\) inner
product. The packed opening protocol authenticates the resulting
challenge-dependent operands against a fixed-basis commitment.
This combination addresses both the relation and its commitment
interface, developed in Sections~\ref{sec:overview} and~\ref{sec:pcs}.
